\capitulo{5}{Aspectos relevantes del desarrollo del proyecto}

Este capítulo tiene como objetivo recoger los aspectos más interesantes y no triviales del proceso de desarrollo del prototipo del laboratorio virtual de programación y robótica. Se detallarán las decisiones de diseño arquitectónico y de implementación clave, justificando los caminos adoptados y las estrategias empleadas para resolver los desafíos técnicos, especialmente aquellos inherentes a la integración de UBlockly en Unity y la simulación del comportamiento de los bloques. Se busca que este apartado sirva como un resumen de la experiencia práctica del proyecto y una referencia útil para futuros desarrolladores.

\section{Metodología de Desarrollo: Iteraciones Ágiles con SCRUM}
El proceso de desarrollo de este prototipo ha seguido una metodología adaptativa, inspirada en los principios de \textbf{SCRUM}. Esta elección, frente a modelos más lineales, se justificó por la naturaleza experimental del proyecto y la necesidad de una \textit{respuesta rápida} a los desafíos tecnológicos emergentes. La aplicación de este marco ágil permitió:
\begin{itemize}
	\item \textbf{Entrega de Valor Incremental:} El proyecto se estructuró en ciclos de \textit{sprints} de duración fija (ej., una a tres semanas), donde al final de cada iteración se obtenía un incremento funcional y probado del prototipo. Esto aseguró un progreso tangible y la validación temprana de funcionalidades críticas, como el movimiento del robot, que era la piedra angular del prototipo.
	\item \textbf{Flexibilidad y Adaptación Continua:} La revisión constante del progreso con el tutor (asumiendo roles de \textit{Product Owner} y \textit{Scrum Master}) permitió reevaluar prioridades, adaptar requisitos y refinar las soluciones técnicas en función de la experiencia práctica obtenida, así como gestionar la complejidad que surgió de la interacción entre UBlockly y Unity.
	\item \textbf{Transparencia y Gestión del Trabajo:} Herramientas como \textbf{Zube} \cite{ZubePlatform} se emplearon para la planificación de \textit{sprints}, el seguimiento de tareas con tableros \textit{Kanban} y la asignación de \textit{puntos de historia}. Esto facilitó la visualización del trabajo pendiente y realizado, mejorando la organización del proyecto.
\end{itemize}
Este enfoque iterativo fue clave para la gestión eficaz de la complejidad técnica del proyecto y para alcanzar un \textit{Producto Mínimo Viable (MVP)} dentro del plazo establecido.

\section{Decisiones de Diseño Arquitectónico e Implementación}

\subsection{Adopción de la Arquitectura Modelo-Vista-Controlador (MVC)}
Una de las decisiones arquitectónicas fundamentales fue la adopción de un patrón \textbf{Modelo-Vista-Controlador (MVC)}. Esta elección es crítica en sistemas complejos como los entornos de desarrollo integrados, donde una separación clara de responsabilidades es esencial para la mantenibilidad, la escalabilidad y la depuración. El MVC se ha implementado de la siguiente manera:
\begin{itemize}
	\item \textbf{Modelo (Lógica de Negocio y Datos Puros):} Las clases centrales del Modelo (\texttt{WorkSpaceModel}, \texttt{BlockModel}, \texttt{ConnectionModel}, \texttt{InputModel}, \texttt{FieldModel}, \texttt{VariableMap}, \texttt{ProcedureDB}) encapsulan el estado lógico y las reglas de negocio del sistema de bloques. Son independientes de cualquier representación visual y se enfocan en la \textit{abstracción de la programación visual} del usuario.
	\item \textbf{Vista (Representación Visual e Interacción UI):} Componentes Unity UI como \texttt{UnityEngine.UI.Image}, \texttt{TMPro.TextMeshProUGUI}, \texttt{UnityEngine.UI.Button} actúan como los elementos visuales, gestionados por una jerarquía de clases \texttt{View} (\texttt{BaseView}, \texttt{BlockView}, \texttt{ConnectionView}, \texttt{InputView}, \texttt{FieldView}, \texttt{LineGroupView}). Estas vistas son las únicas que conocen la interfaz de Unity y se encargan de representar el Modelo.
	\item \textbf{Controlador (Orquestación e Interpretación de Interacciones):} Capas de lógica como \texttt{AppController}, \texttt{WorkspaceController}, \texttt{BlockDragController}, \texttt{BlockConnectionController}, \texttt{ExecutionController} e \texttt{InputController} median entre el Modelo y la Vista. Interpretan las acciones del usuario, validan operaciones y orquestan las actualizaciones del Modelo y la Vista sin que estas capas se conozcan directamente entre sí.
\end{itemize}
La comunicación entre estas capas se gestiona primordialmente mediante el \textbf{patrón Observador} (\texttt{Observable<TArgs>} e \texttt{IObserver<TArgs>}). Los modelos (\texttt{BlockModel}, \texttt{ConnectionModel}, \texttt{FieldModel}) actúan como sujetos observados, notificando de cualquier cambio relevante a sus vistas o controladores suscritos. Esto asegura una \textit{sincronización eficiente} de la interfaz sin acoplamientos rígidos entre los elementos. Por ejemplo, cuando un usuario cambia el valor de un campo (\texttt{FieldView}), este notifica al controlador (\texttt{InputController}), que solicita al modelo (\texttt{FieldModel}) el cambio, el cual a su vez notifica a la vista (\texttt{FieldView}) para que se actualice.

\subsection{Aprovechamiento y Adaptación del Repositorio UBlockly}
La decisión de basar la implementación del motor de bloques en el repositorio de \textbf{UBlockly} \cite{UBlockly}, una adaptación de Blockly en C\# para Unity, fue una \textit{elección estratégica} frente al desarrollo desde cero. Este enfoque se justificó por la necesidad de construir un prototipo funcional robusto en un tiempo limitado, ya que UBlockly proporcionó una base sólida para:
\begin{itemize}
	\item \textbf{Motor Lógico de Bloques:} Sus clases core como \texttt{WorkSpaceModel} (gestor principal de la base de datos de bloques), \texttt{BlockModel} (estructura lógica de un bloque) y \texttt{ConnectionModel} (sistema de conexiones entre bloques) ofrecieron una base abstracta y funcional ya testeada. Se adoptó este diseño tal cual, centrándose la adaptación en cómo estos modelos interactúan con la nueva capa de vistas en Unity.
	\item \textbf{Motor de Ejecución Asíncrono:} UBlockly ya incluye un robusto sistema de interpretación directa basado en \texttt{Cmdtor}s y corrutinas de C\# (\texttt{CSharpRunner}). Esta arquitectura se ha adoptado directamente, adaptando los intérpretes (\texttt{MotionBlockInterpreter}) para interactuar con los objetos simulados del robot.
	\item \textbf{Serialización y Persistencia:} El mecanismo de UBlockly para serializar el estado del \texttt{Workspace} a \texttt{XML} y viceversa se integró directamente, proporcionando una solución eficiente para guardar y cargar los proyectos del usuario.
\end{itemize}
La adaptación principal ha consistido en \textit{acoplar las capas de modelo y ejecución de UBlockly con la nueva capa de visualización de Unity UI}, así como \textit{customizar la lógica de interacción} (drag \& drop) para lograr una experiencia más cercana a Scratch. Se han creado prefabs de GameObjects de Unity con la estructura de componentes \texttt{View} que reflejan los \texttt{BlockModel}s lógicos, permitiendo la creación modular y una apariencia flexible de los bloques.

\section{Detalles Clave de Implementación}

\subsection{Representación Visual y Composición de Bloques en Unity}
Uno de los mayores desafíos en el desarrollo de interfaces de programación visual es la correcta representación y sincronización entre el modelo lógico del programa y su visualización interactiva. En este proyecto, la aproximación se basa en la creación de \textbf{prefabs} de \texttt{GameObject}s de Unity para cada tipo de bloque (\texttt{BlockView}), replicando la estética de Scratch.
\begin{itemize}
	\item \textbf{Creación de Prefabs Estructurados:} Cada \texttt{BlockView} es un prefab con una jerarquía de componentes \texttt{UnityEngine.UI} que se corresponden directamente con las estructuras lógicas del \texttt{BlockModel} (ej., \texttt{InputView} para entradas, \texttt{FieldView} para campos). Por ejemplo, el bloque \texttt{motion\_movesteps} es un prefab con un \texttt{GameObject} \texttt{motion\_movesteps} (con \texttt{BlockView} y \texttt{Image}) que agrupa lógicamente los hijos (\texttt{LineGroup}, \texttt{Connection\_prev}, \texttt{Connection\_next}, \texttt{nextStatementContainer}). Esta composición predefinida de prefabs acelera la instanciación de bloques en tiempo de ejecución.
	\item \textbf{Layout Manual y Sincronización de \texttt{RectTransform}:} A diferencia de los sistemas de layout automático de Unity, el proyecto emplea un \textbf{sistema de layout manual} (\texttt{BaseView.UpdateLayout()}) para los bloques. Este sistema calcula dinámicamente la posición (\texttt{XY}) y el tamaño (\texttt{CalculateSize()}) de cada \texttt{BlockView} y sus subcomponentes (\texttt{InputView}, \texttt{FieldView}, \texttt{ConnectionView}). Es crucial porque los bloques de programación visual tienen formas complejas (muescas, pestañas, huecos para entradas) que no se ajustan de manera óptima a layouts grid o verticales/horizontales simples. Cada vez que el modelo lógico de un bloque cambia de posición (\texttt{BlockModel.XY}), o su contenido o tamaño (\texttt{BlockModel.FireUpdate()} para \texttt{UpdateStates.Inputs} o \texttt{UpdateStates.Connections}), el \texttt{BlockView} es notificado y marca su layout como sucio (`NotifyLayoutDirty`). Posteriormente, en un ciclo de Unity, se fuerza un recálculo del layout completo o parcial mediante \texttt{LayoutRebuilder.ForceRebuildLayoutImmediate()}.
	\item \textbf{Visualización de Conexiones e Identificación de \texttt{nextStatementContainer}:} Las conexiones lógicas (\texttt{ConnectionModel}) tienen su reflejo visual en las \texttt{ConnectionView}s. Un elemento fundamental en el diseño de los prefabs es el \texttt{RectTransform} denominado \texttt{nextStatementContainer} (ver jerarquía en la imagen adjunta en Capítulo 3 para \texttt{motion\_movesteps}). Este \texttt{GameObject} no solo es un punto de anclaje visual, sino el \textbf{contenedor principal para la pila de bloques subsiguientes} conectados al \texttt{NextConnection}. Cuando un bloque se une, su \texttt{BlockView} raíz se re-parenta a este \texttt{nextStatementContainer}, asegurando la correcta representación jerárquica y el apilamiento visual de las sentencias.
\end{itemize}

\subsection{Interacción con los Bloques: El Sistema de Arrastre y Conexión (Drag \& Drop)}
La funcionalidad de arrastrar y conectar bloques de forma intuitiva y fluida es central para replicar la experiencia de Scratch. Este sistema, orquestado por \texttt{BlockDragController} y \texttt{BlockConnectionController}, representa una implementación compleja con diversas interacciones:
\begin{itemize}
	\item \textbf{Inicio del Arrastre:} Se gestiona desde un \texttt{BlockTemplateDragSource} para las plantillas de la \texttt{Toolbox}, o directamente desde la \texttt{BlockView} para bloques ya en el área de codificación. Este proceso implica la clonación del \texttt{BlockModel} (si es una plantilla), la creación de su \texttt{BlockView} asociada y su re-parentado a una capa superior (`DragLayer`) para la interacción flotante. Se guarda un offset ('m\_dragStartOffset`) para mantener la posición del puntero relativa al bloque.
	\item \textbf{Búsqueda y Resaltado de Conexiones:} Durante el arrastre, el \texttt{BlockDragController} solicita a \texttt{BlockConnectionController.ProcessDrag()} la identificación de posibles conexiones válidas. Este proceso implica:
	\begin{itemize}
		\item Recorrer todas las \texttt{ConnectionModel}s del bloque que se está arrastrando.
		\item Utilizar las \texttt{BlockConnectionDB}s (\texttt{ConnectionDBList} en el \texttt{WorkSpaceModel}) para buscar eficientemente las conexiones compatibles más cercanas (\texttt{SearchForClosest()}) a los puntos de conexión del bloque arrastrado. Estas bases de datos en memoria están optimizadas espacialmente para la búsqueda de vecinos cercanos, manteniendo la \texttt{Location} (\texttt{Vector2}) actualizada para cada \texttt{ConnectionModel}.
		\item Aplicar reglas de validación estricta (\texttt{IsConnectionAllowed()}) para determinar si una conexión es posible, incluyendo compatibilidad de tipos, prevención de bucles jerárquicos y distancia límite (snap).
		\item Si se encuentra un candidato, la \texttt{ConnectionView} correspondiente es instruida para activar un \textit{resaltado visual} (`Highlight`), indicando al usuario la posible unión.
	\end{itemize}
	\item \textbf{Conexión Lógica y Re-parentado Visual (\texttt{Snap}):} El punto más crítico se da al soltar el bloque (\texttt{BlockDragController.HandleEndDrag()}). Si existe una conexión válida:
	\begin{itemize}
		\item \textbf{Conexión Lógica:} Las dos \texttt{ConnectionModel}s (del bloque arrastrado y del estático) se unen mediante \texttt{ConnectionModel.Connect()}. Este método es fundamental por su implementación de la \textit{lógica de bumping}: si una conexión receptora ya está ocupada, el bloque previamente conectado se \textit{desconecta y se desplaza} (`UpdateState.BumpedAway`) a una posición cercana, y se busca una nueva ubicación para él dentro de la jerarquía de bloques (`ConnectionModel.LastConnectionInRow()`).
		\item \textbf{Re-parentado Visual:} La \texttt{BlockView} del bloque conectado se re-parenta inmediatamente al \texttt{RectTransform} adecuado de su nuevo padre visual (ej. \texttt{nextStatementContainer}), ajustando su posición \texttt{anchoredPosition} para que encaje visualmente.
		\item \textbf{Actualización de Layouts:} Tras cualquier conexión o desconexión, el sistema marca los \texttt{RectTransform}s afectados para una \textit{reconstrucción del layout}, asegurando que la interfaz se adapte visualmente al nuevo estado del programa (\texttt{LayoutRebuilder.MarkLayoutForRebuild()}).
	\end{itemize}
\end{itemize}

\subsection{El Motor de Ejecución de Bloques: De la Lógica a la Acción del Robot}
El corazón operativo del prototipo es el motor de ejecución de bloques, responsable de interpretar la secuencia de instrucciones del usuario y traducirlas en acciones en el entorno de simulación 3D.
\begin{itemize}
	\item \textbf{Intérpretes de Bloques (\texttt{Cmdtor}s):} Cada tipo de bloque funcional (ej. movimiento) tiene una clase intérprete concreta que hereda de \texttt{Cmdtor} (\texttt{ValueCmdtor}, \texttt{VoidCmdtor}, \texttt{EnumeratorCmdtor}). Por ejemplo, \texttt{MotionBlockInterpreter} se encarga de extraer los parámetros de mover N pasos de un \texttt{BlockModel}.
	\item \textbf{Orquestación con \texttt{CSharpRunner}:} La clase \texttt{CSharpRunner}, un \texttt{MonoBehaviour}, es el orquestador principal de la ejecución. Mantiene una \texttt{Stack<IEnumerator>} de \texttt{CmdEnumerator}s para gestionar el flujo de control (secuencia lineal, bucles, condicionales). Esto permite una \textit{ejecución paso a paso no bloqueante} del programa.
	\item \textbf{Concurrencia con Corrutinas:} La característica distintiva de este sistema es el uso intensivo de las \textbf{corruutinas de C\# (\texttt{IEnumerator}, \texttt{yield return})}. Cuando un intérprete de bloque (\texttt{Cmdtor}) invoca una acción que toma tiempo (ej., una animación de movimiento del robot, un retardo), este se traduce en un \texttt{yield return} en C\#. El \texttt{CSharpRunner} \textit{pausa la ejecución de ese bloque} y la \textit{reanuda automáticamente} cuando la corrutina de la acción externa ha finalizado. Esto es crucial para lograr \textit{animaciones fluidas} del robot o \textit{retrasos precisos} sin congelar la aplicación. Por ejemplo, \texttt{MotionBlockInterpreter} llama a \texttt{BlockObserver.\_MoveRobot()} (una corrutina), y el \texttt{CSharpRunner} espera a que esta corrutina termine antes de pasar al siguiente bloque.
	\item \textbf{Sincronización con la Simulación:} La clase \texttt{BlockObserver} actúa como un puente entre el motor de ejecución de bloques y el \texttt{GameObject} del robot simulado (\texttt{RobotBehaviour}). \texttt{BlockObserver} recibe las órdenes de alto nivel (ej. mover 10 pasos) desde los \texttt{Cmdtor}s y delega estas acciones a \texttt{RobotBehaviour} (el script que controla el movimiento y la animación física del robot), asegurando que la simulación refleje con precisión la lógica del programa. \texttt{BlockObserver} también gestiona el feedback visual en la interfaz de usuario, como los mensajes de estado y el resaltado del bloque en ejecución.
\end{itemize}

\subsection{Persistencia de Proyectos del Usuario}
La capacidad de guardar y cargar los programas creados por el usuario es esencial para la continuidad del aprendizaje y la gestión de proyectos. El prototipo implementa un sistema básico pero funcional de persistencia.
\begin{itemize}
	\item \textbf{Serialización en Formato XML:} La estructura lógica completa del \texttt{WorkSpaceModel} (incluyendo todos los bloques, sus jerarquías internas y las conexiones establecidas) se serializa a un formato \textbf{XML}. Las clases \texttt{Xml} de UBlockly (\texttt{WorkspaceToDom()} y \texttt{DomToWorkspace()}) son el pilar de este proceso, transformando el árbol de objetos C\# del modelo a un árbol XML y viceversa. Esta elección del formato XML proporciona una representación clara y legible de los programas.
	\item \textbf{Almacenamiento Local con \texttt{PlayerPrefs}:} La clase \texttt{WorkspaceSerializer}, desarrollada para este proyecto, es la encargada de manejar la persistencia del \texttt{XML} generado. Opta por \textit{guardar el XML resultante directamente en las \texttt{PlayerPrefs}} de Unity y \textit{cargarlo} desde esta ubicación local. Esta solución, aunque simple, es altamente eficaz para un prototipo, dada su facilidad de implementación y su funcionalidad de almacenamiento rápido para la continuidad en los proyectos de usuario en un único dispositivo. Es importante destacar que, para este prototipo, se ha priorizado esta solución frente a alternativas más robustas o escalables como bases de datos en tiempo real (ej. Firebase), las cuales se considerarían para fases de desarrollo futuras.
\end{itemize}

\section{Retos Técnicos Abordados y Estrategias de Depuración}
El desarrollo de un sistema híbrido de esta complejidad, integrando una librería de lógica de bloques con el entorno 3D y UI de Unity, ha presentado desafíos técnicos significativos.
\begin{itemize}
	\item \textbf{Sincronización de Modelos y Vistas en MVC con Flotantes y Corrutinas:} El reto de mantener el \texttt{BlockModel} lógico (con posiciones \texttt{Vector2} abstractas) perfectamente sincronizado con el \texttt{RectTransform} visual en el \texttt{Unity UI Canvas} ha requerido una depuración exhaustiva. Problemas de desajustes visuales o saltos en la posición de los bloques (especialmente al soltarlos tras un arrastre o al apilarse), así como errores de estado (`BlockView` o `ConnectionView` mostrando un estado inconsistente con el modelo), han sido frecuentes. La complejidad aumenta con la naturaleza flotante de las coordenadas y la ejecución asíncrona de las animaciones del robot, donde la interfaz debe reflejar cambios precisos en momentos concretos.
	\item \textbf{Bases de Datos de Conexiones en Memoria (\texttt{BlockConnectionDB}):} Las \texttt{BlockConnectionDB}s de UBlockly utilizan estructuras de datos como cuadrantes o segmentos espaciales para optimizar la búsqueda de conexiones cercanas durante el arrastre. Asegurar que cada \texttt{ConnectionModel} estuviera correctamente registrado y desregistrado de estas bases de datos (`ConnectionModel.Location` debía estar perfectamente actualizada, y las funciones `ConnectionModel.AddConnection()` y `ConnectionModel.RemoveConnection()` sin errores de referencia a objetos ya destruidos) ha sido un punto crítico. La coherencia de estas estructuras internas, que son abstractas (en memoria) y no se ven directamente en la UI, es vital para que el sistema de snap funcione.
	\item \textbf{Gestión de la Jerarquía de \texttt{GameObject}s y Disposición (\texttt{Dispose}/\texttt{UnPlug}):} La correcta re-parentalización de \texttt{GameObject}s \texttt{BlockView} al conectarse y desconectarse (particularmente la lógica de bumping) y la eliminación de recursos ha requerido un cuidado extremo. Un error aquí podía llevar a fugas de memoria, bloques invisibles, bloques duplicados o zombies en la jerarquía, o que el robot saliera lanzado por una lógica de corrutinas mal coordinada que no espera la finalización de las animaciones.
	
	\item \textbf{Estrategias de Depuración para Sistemas Complejos y Asíncronos:} Dada la complejidad y la interacción entre múltiples componentes MVC, corrutinas y la interfaz de Unity, las herramientas de depuración han sido esenciales:
	\begin{itemize}
		\item \textbf{Logueo Exhaustivo y Contextualizado (\texttt{Logger}):} Se ha implementado y utilizado de forma intensiva la clase \texttt{Logger}, una utilidad global para la emisión de mensajes `Debug.Log`, `Debug.LogWarning` y `Debug.LogError`. \textit{Decenas de puntos críticos en el código incluyen `Debug.Log`s detallados (a menudo comentados para evitar el ruido en la ejecución normal) que rastrean el estado exacto de variables, la entrada y salida de funciones, el estado de las corrutinas, y la pertenencia de objetos a bases de datos.} Estos logs han sido el principal mecanismo para seguir el flujo de control, identificar desajustes y diagnosticar problemas de temporización.
		\item \textbf{Herramientas Visuales de Depuración en Unity:} El editor de Unity ha permitido la inspección en tiempo real de la jerarquía de \texttt{GameObject}s, los valores de los \texttt{RectTransform}s, el estado de los componentes \texttt{BlockView} y \texttt{ConnectionView}. Las herramientas de \textbf{Debugging de la ConnectionDB (\texttt{ExportConnectionDBsButton} y \texttt{WorkspaceController.ExportConnectionDBState()})}, desarrolladas para el prototipo, han sido fundamentales para \textit{visualizar y exportar a JSON} el estado de las bases de datos de conexiones en memoria, permitiendo un análisis detallado de su coherencia. Esto ha sido invaluable para la depuración de la lógica de arrastre y snap.
	\end{itemize}
	La combinación de estas estrategias ha sido fundamental para navegar por la complejidad del desarrollo, asegurar la estabilidad del prototipo y justificar las decisiones de diseño a través de la evidencia empírica generada en el proceso de depuración.
\end{itemize}